\documentclass[11pt,a4paper]{article}
\usepackage[T1]{fontenc}
\usepackage[utf8]{inputenc}
\usepackage{amsmath,amssymb,amsthm}
\usepackage[margin=1in]{geometry}
\usepackage[colorlinks=true,linkcolor=blue,urlcolor=blue]{hyperref}
\theoremstyle{plain}
\newtheorem{theorem}{Theorem}
\newtheorem{lemma}[theorem]{Lemma}
\newtheorem{proposition}[theorem]{Proposition}
\newtheorem{corollary}[theorem]{Corollary}
\theoremstyle{definition}
\newtheorem{remark}[theorem]{Remark}

\title{The empirical recurrence for OEIS A257359, proved by transfer matrix}
\author{Adrian Perez Fontelles\\ \small Independent researcher}
\date{3 September 2026}
\begin{document}
\maketitle

\begin{abstract}
OEIS A257359 counts arrays of width $8$ over $\{0,1\}$ in which every $3\times3$ window has
constrained SUMS along its rows, its columns and its two diagonals. The entry carries an
empirical recurrence of order $42$ contributed by R.~H.~Hardin. It is true, and it is
decidable rather than empirical: a $3\times3$ window lies in $3$ consecutive array rows,
so the count is a walk count on $S=65793$ states.
\end{abstract}

\noindent\small 2020 Mathematics Subject Classification. 05A15, 05B20, 15A18.\normalsize

\section{The conjecture}

OEIS A257359 is ``Number of (n+2)X(6+2) 0..1 arrays with no 3x3 subblock diagonal sum equal to the antidiagonal sum or central row sum less than the central column sum.''. It has offset $1$ and begins
\[
174724, 145924, 234256, 622521, 1196836, 2433600, 6012304, 15547249,\ \dots
\]
The entry states, as a conjecture contributed by its author and never marked settled:
\begin{quote}
Empirical: a(n) = 4*a(n-2) +17*a(n-4) -57*a(n-6) +2*a(n-8) -60*a(n-10) -127*a(n-12) +598*a(n-14) -292*a(n-16) +315*a(n-18) +312*a(n-20) -1149*a(n-22) +223*a(n-26) -55*a(n-28) +127*a(n-30) +51*a(n-32) -32*a(n-34) -6*a(n-36) -5*a(n-38) -a(n-40) +a(n-42) for n>48
\end{quote}
As of the ``Last modified'' line on the live entry (Jul 23 2025 15:22:00, revision 6) this is still
recorded as empirical, and nothing on the entry records it as proved.

\section{What is being counted}

The array has $8$ cells across, with entries in $\{0,1\}$. A WINDOW is a $3\times3$ block
of consecutive cells lying inside the array. Every window must satisfy
\[
\begin{array}{ll}
\text{the sum along} & \text{} \\[2pt]
\text{the main diagonal} & \text{must NOT be equal to the sum along the antidiagonal} \\
\text{the central row} & \text{must NOT be less than the sum along the central column} \\
\end{array}
\]
where the main diagonal of a window runs from its top left corner to its bottom right, and its
antidiagonal from top right to bottom left.

\section{The walk}

\begin{lemma}\label{lem:win}
A $3\times3$ window occupies $3$ consecutive array rows, so every condition above is
decided by $3$ consecutive rows and by nothing else.
\end{lemma}

\begin{proof}
By definition of a window.
\end{proof}

So the state is the window of the last $2$ array rows, a row not yet written carried as
$\bot$, and a step appends a row and settles every $3\times3$ window that row completes. A
window needing a row that is $\bot$ does not exist, which is the boundary rule; the walk starts
at the all-$\bot$ state, so a walk of $R$ steps is an array of $R$ rows. With $M$ the adjacency
matrix and $\tau=\mathbf{1}$,
\[
a(n)\;=\;\iota^{\!\top}M^{\,j}\tau ,\qquad S=65793.
\]
The walk index $j$ and the entry's own $n$ differ by a constant, read off by matching the
published terms rather than assumed. In particular $a$ is $C$-finite.

\section{The proof}

\begin{lemma}\label{lem:crit}
Let $q(t)=t^{42}-\sum_i c_it^{42-i}$ be the characteristic polynomial of the
conjectured recurrence. The residual $a(n)-\sum_ic_ia(n-i)$ equals
$\iota^{\!\top}M^{\,j}q(M)\tau$ for the corresponding walk index $j$, so the recurrence
holds from some point on if and only if $\iota^{\!\top}M^{j}q(M)\tau=0$ for all large $j$,
and it is enough to examine $j<S$.
\end{lemma}

\begin{proof}
Factor $M^{j}$ out of $M^{j+42}-\sum_ic_iM^{j+42-i}$. For the bound, $M^{S}$
is an integer combination of $I,M,\dots,M^{S-1}$ by Cayley--Hamilton, so the numbers
$u_j=\iota^{\!\top}M^{j}q(M)\tau$ obey the monic recurrence given by the characteristic
polynomial of $M$; $S$ consecutive zeros force every later one, and the last nonzero $u_j$
pins the threshold exactly.
\end{proof}

\begin{theorem}
\[
a(n)\;=\;4\,a(n-2) + 17\,a(n-4) - 57\,a(n-6) + 2\,a(n-8) - 60\,a(n-10) - 127\,a(n-12) + 598\,a(n-14) - 292\,a(n-16) + 315\,a(n-18) + 312\,a(n-20) - 1149\,a(n-22) + 223\,a(n-26) - 55\,a(n-28) + 127\,a(n-30) + 51\,a(n-32) - 32\,a(n-34) - 6\,a(n-36) - 5\,a(n-38) - a(n-40) + a(n-42)
\]
for every $n>48$.
\end{theorem}

\begin{proof}
The vector $w=q(M)\tau$ was formed by $42$ matrix--vector products in exact integer
arithmetic and $\iota^{\!\top}M^{j}w$ evaluated for $j=0,1,\dots$, again exactly; those
integers vanish from the index corresponding to $n=48+1$ onwards and the run continued
until the working vector was identically zero.
\end{proof}

\section{Verification}

Four checks, each independent of the algebra above.

First, the model reproduces all $21$ terms the entry publishes, exactly, in integer
arithmetic.

Second, the count was recomputed for the smallest arrays straight from the definition: fill the
cells in row major order, in the orientation the entry's own name writes, and test a window as
soon as its last cell is placed. No state, no transposition and no transfer matrix are
involved, and the published terms came back.

Third, the conjectured recurrence was evaluated directly on the published terms, with no
matrices involved, and holds wherever the proved range and the data overlap.

Fourth, the annihilation test of Lemma~\ref{lem:crit} was carried out in exact integer
arithmetic on the whole vector rather than on a sampled prefix.

\begin{thebibliography}{9}
\bibitem{oeis} The OEIS Foundation, \emph{The On-Line Encyclopedia of Integer
Sequences}, \url{https://oeis.org}, 2026. Sequence A257359.
\bibitem{stanley} R.~P.~Stanley, \emph{Enumerative Combinatorics}, Volume 1, 2nd ed.,
Cambridge University Press, 2011. (Transfer-matrix method, Section 4.7.)
\bibitem{horn} R.~A.~Horn and C.~R.~Johnson, \emph{Matrix Analysis}, 2nd ed., Cambridge
University Press, 2013.
\end{thebibliography}

\end{document}
